\chapter{Worked examples}
\label{ch:examples}

The repository carries eighteen runnable examples. This chapter is a guide to
which one answers which question, followed by three walked through in enough
detail to copy from.

\section{The map}

\begin{center}
\small
\begin{longtable}{@{}L{5.1cm}L{8.5cm}@{}}
\toprule
\textbf{Path} & \textbf{What it shows} \\
\midrule
\endhead
\code{examples/manifest} &
  Manifest-driven generation for \code{Person}, with no class modification. The
  smallest complete case. \\
\addlinespace
\code{examples/annotate-manifest} &
  Out-of-line annotations from an external JSON file, wired by the manifest's
  \field{annotations} field. \\
\addlinespace
\code{examples/doxygen} &
  The documentation a library already has: plain Doxygen comments and default
  arguments in an untouched header become Python docstrings and keyword
  arguments, TSDoc, OpenAPI descriptions and a Markdown reference.
  Ch.~\ref{ch:documentation} \\
\addlinespace
\code{examples/geom-lib} &
  Manifest-driven bindings for a small geometry library --- nested types,
  vectors. \\
\addlinespace
\code{examples/geom-expanded} &
  Ten backends (\lang{python}, \lang{nanobind}, \lang{node}, \lang{wasm},
  \lang{qt}, \lang{qml}, \lang{csharp}, \lang{java}, \lang{lua}, \lang{julia})
  built with off-the-shelf toolchains, plus out-of-line annotations. The proof
  of Section~\ref{sec:reflection-free}. \\
\addlinespace
\code{examples/dynamic} &
  The \lang{dynamic} backend end to end: one set of generated metadata driving a
  terminal interpreter \emph{and} a Qt viewer (3-D view, property panel,
  console), neither naming a bound type. \\
\addlinespace
\code{examples/scriptable-model} &
  The two-stage pipeline of Chapter~\ref{ch:scriptable}. \\
\addlinespace
\code{examples/plain-binding} &
  One library, one manifest, four language bindings --- the base case, run
  once. \\
\addlinespace
\code{examples/trampoline-python} &
  Overriding C++ virtuals from Python: generated pybind11 trampolines from
  \code{virtual\_spec}. \\
\addlinespace
\code{examples/trampoline-node} &
  The same from JavaScript, with N-API trampolines. \\
\addlinespace
\code{examples/dynamic-lib} &
  Linking the bindings against a pre-built library --- the \field{user\_lib}
  case. \\
\addlinespace
\code{examples/moc} &
  The moc-less signals/slots/properties demo built on \code{mini\_moc.h}. \\
\addlinespace
\code{examples/docgen} &
  A reflection-driven Markdown / HTML reference generator. \\
\addlinespace
\code{examples/paraview} &
  ParaView plugin property-panel XML from an annotated \code{vtkThreshold}
  spec --- exercises every annotation feature. \\
\addlinespace
\code{examples/qt}, \code{examples/qml}, \code{examples/imgui} &
  The three inspector backends, each building a form from a reflected struct.
  \code{imgui} uses an \code{Algo.ann.json} side-car and builds with an
  off-the-shelf C++20 compiler, dependencies auto-fetched. \\
\addlinespace
\code{examples/bindings/*} &
  Hand-written reference backends --- pybind11, N-API, jlcxx, REST, embind.
  What the generator emits, written by hand, for comparison. \\
\bottomrule
\end{longtable}
\end{center}

\section{The smallest case}

\code{examples/manifest} is the one to read first. An unmodified
\code{person.h}, a manifest of ten lines, and the pipeline. It is
Chapter~\ref{ch:first} as a runnable directory.

\section{A geometry library}

\code{examples/geom-lib} is the first realistic shape: several classes, one
returning a \code{std::vector} of another, an enum, and a free function.

The two things to look at:

\begin{itemize}
  \item \textbf{Nested user types cross without being mentioned.} A
    \code{Surface} method returning \code{std::vector<Triangle>} works because
    \code{Triangle} is in \field{classes} --- not because anything says how to
    convert it.
  \item \textbf{One module per target, not one per class.} The generated
    \code{pygeom} carries \code{Model}, \code{Point}, \code{Surface},
    \code{Triangle} and \code{Kind} together.
\end{itemize}

\code{examples/geom-expanded} takes the same library and builds ten backends
with stock toolchains --- no fork anywhere in the second pass. If you want to
convince yourself (or a colleague) that the reflection-free claim is real, that
is the directory to run.

\section{The dynamic model, end to end}

\code{examples/dynamic} is the most interesting example in the repository and
the one that takes longest to appreciate.

A small scene library --- \code{Vec3}, \code{Mesh}, and a \code{Relaxer} solver
--- goes through the \lang{dynamic} backend once. Two entirely separate
consumers then drive it:

\begin{itemize}
  \item a \textbf{terminal interpreter} that lists types, gets and sets fields,
    and calls methods by name;
  \item a \textbf{Qt viewer} with a 3-D view, a property panel and a console.
\end{itemize}

Neither names a single bound type. The property panel is a walk over
\code{fields()}: \code{doc} becomes a tooltip, \code{range} becomes slider
bounds, \code{combobox} becomes a drop-down.

\subsection*{What adding a class costs}

The example's own history is the argument. Adding \code{scene::Relaxer} ---
tangential Laplacian relaxation, a genuine algorithm with tuning parameters and
one method that runs for a while --- touched three files, \emph{all of them the
library's own}:

\begin{enumerate}
  \item \code{scene.h}, to add the class;
  \item \code{Mesh.ann.json}, for one new accessor's annotation;
  \item one line of \code{manifest.json}.
\end{enumerate}

No consumer changed. Not the Qt viewer, not the interpreter, not any of the five
scriptable drivers --- yet all of them gained it: the solver appeared in the
type list, its parameters appeared in the property panel with their declared
ranges, and its method appeared as a button.

That is what ``UI by query rather than by codegen'' buys, stated as a diff.

\section{Rosetta in the wild}

Real libraries bound with rosetta, each from a single manifest and with no
hand-written wrappers:

\begin{center}
\small
\begin{tabular}{@{}L{3.2cm}L{5.0cm}L{4.6cm}@{}}
\toprule
\textbf{Project} & \textbf{Bound library} & \textbf{Targets} \\
\midrule
\code{pmp-rosetta} & PMP --- polygon mesh processing (remeshing, smoothing,
  subdivision, decimation) & Python, Node.js, WebAssembly, TypeScript \\
\addlinespace
\code{geogram-rosetta} & geogram --- geometry processing (reconstruction,
  remeshing, parameterisation, booleans) & Python, Node.js, WebAssembly,
  TypeScript, Lua \\
\addlinespace
\code{arch-rosetta} & Arch --- 3-D boundary-element geomechanics &
  Python, Node.js, WebAssembly, TypeScript \\
\addlinespace
\code{cassini-rosetta} & Cassini --- FEM geomechanical restoration, through a
  single high-level facade & Python, Node.js, WebAssembly, TypeScript \\
\addlinespace
\code{triax-rosetta} & Triax --- discrete-element triaxial compression &
  Python, Node.js, WebAssembly \\
\bottomrule
\end{tabular}
\end{center}

The geogram manifest is the best worked reference for a difficult library: it
uses \field{sequences} for \code{GEO::vector}, \field{extensions} to reach
geometry behind raw-pointer accessors, \field{signature} to pick one overload of
\code{mesh\_union}, \field{module\_init} for \code{GEO::initialize()},
\field{generated\_headers} for the configured \code{version.h}, and
\field{link\_options} for the wasm NODEFS mount --- in other words, nearly every
field in Part~II, each for a reason you can see.

\section{Extending a generated binding by hand}

Everything under \code{bindings/} is regenerated output and must never be
edited. When the generated binding misses something you need --- a helper the
walker skips, a typed-array export for a renderer --- add it in a
\textbf{separate hand-written file} compiled \emph{alongside} the generated
source, from a small build of your own that lives outside \code{bindings/}.

The binding frameworks accept several registration blocks per module, so your
file contributes a second one. Nothing generated is touched, and regenerating
never clobbers your extensions.

\begin{lstlisting}[style=cpp]
// viz_helpers.cpp -- compiled next to the generated auto_emscripten.cpp
emscripten::val mesh_positions(const pmp::SurfaceMesh &mesh);  // Float32Array
emscripten::val mesh_triangles(const pmp::SurfaceMesh &mesh);  // Uint32Array
void triangulate_mesh(pmp::SurfaceMesh &m) { pmp::triangulate(m); }

EMSCRIPTEN_BINDINGS(pmp_viz) {   // a second block; the generated one keeps its own
    emscripten::function("mesh_positions", &mesh_positions);
    emscripten::function("mesh_triangles", &mesh_triangles);
    emscripten::function("triangulate_mesh", &triangulate_mesh);
}
\end{lstlisting}

\begin{warn}
Embind is the friendliest here because it accepts any number of
\code{EMSCRIPTEN\_BINDINGS} blocks in one module. Backends with a single module
entry point --- pybind11's \code{PYBIND11\_MODULE}, N-API's
\code{NODE\_API\_MODULE} --- cannot add a second block to the \emph{same}
module. There, ship your extras as a small companion module built the same way
(it sees the same user headers and links the same library) and import both.
\end{warn}

Before reaching for this, check whether \field{extensions}
(Section~\ref{sec:extensions}) covers your case: a free function taking
\code{Cls\&} needs no build of your own and reaches every backend at once.
